% !TEX root = ../Main.tex
\chapter{韦达定理与弦长}

一元二次方程是初升高衔接最"硬"的一块。韦达定理不仅告诉你根与系数的关系，更是解析几何里"直线与圆锥曲线相交"问题的万能工具——因为\textbf{几乎所有"弦"的信息都能转化为韦达公式的线性组合}。

\section{韦达定理}

\thm{韦达定理}{
    若一元二次方程 $ax^2 + bx + c = 0$（$a \ne 0$）的两根为 $x_1, x_2$，则
    \[
    x_1 + x_2 = -\frac{b}{a},\qquad x_1 x_2 = \frac{c}{a}.
    \]
}

\pf{
    由求根公式 $x_{1,2} = \dfrac{-b \pm \sqrt{b^2 - 4ac}}{2a}$，直接相加相乘即得。

    更简洁的代数证明：$ax^2 + bx + c = a(x - x_1)(x - x_2) = a\bigl[x^2 - (x_1+x_2)x + x_1 x_2\bigr]$，比较系数即得。
}

\rmkb{
    韦达定理的威力在于：很多问题不需要解出 $x_1, x_2$ 各自的值，只需要 $x_1 + x_2$ 和 $x_1 x_2$ 就能推出答案。
}

\section{弦长公式}

设直线 $\ell: y = kx + m$ 与曲线交于 $A(x_1, y_1), B(x_2, y_2)$。则
\[
|AB|^2 = (x_1 - x_2)^2 + (y_1 - y_2)^2 = (x_1 - x_2)^2 (1 + k^2).
\]
又
\[
(x_1 - x_2)^2 = (x_1 + x_2)^2 - 4 x_1 x_2,
\]
故

\state{弦长公式}{
    \[
    |AB| = \sqrt{1 + k^2} \cdot |x_1 - x_2| = \sqrt{1 + k^2} \cdot \sqrt{(x_1 + x_2)^2 - 4 x_1 x_2}.
    \]
    把"交点的横坐标"通过韦达定理整体刻画，避免求具体的交点坐标。
}

\ex[抛物线弦长]{
    直线 $y = x - 1$ 与抛物线 $y^2 = 4x$ 交于 $A, B$。求 $|AB|$。
}

\sol{
    代入得 $(x-1)^2 = 4x$，即 $x^2 - 6x + 1 = 0$。由韦达定理 $x_1 + x_2 = 6$，$x_1 x_2 = 1$。又 $k = 1$。
    \[
    |AB| = \sqrt{1 + 1^2} \cdot \sqrt{6^2 - 4 \cdot 1} = \sqrt{2} \cdot \sqrt{32} = \sqrt{64} = 8.
    \]
}

\ex[椭圆内接三角形的面积]{
    直线 $\ell: y = kx + 1$ 与椭圆 $\dfrac{x^2}{4} + y^2 = 1$ 交于 $A, B$。求 $k$ 使得 $\triangle OAB$ 面积最大（$O$ 为原点）。
}

\sol{
    代入椭圆方程：
    \[
    \frac{x^2}{4} + (kx+1)^2 = 1
    \implies (1 + 4k^2) x^2 + 8kx + 0 = 0.
    \]
    等等，展开仔细：$\tfrac{x^2}{4} + k^2 x^2 + 2kx + 1 = 1$，即 $\left(\tfrac{1}{4} + k^2\right) x^2 + 2kx = 0$，
    \[
    (1 + 4k^2) x^2 + 8kx = 0.
    \]
    由韦达 $x_1 + x_2 = -\tfrac{8k}{1+4k^2}$，$x_1 x_2 = 0$。于是一根为 $0$（恰对应 $\ell$ 与 $y$ 轴的交点 $(0, 1)$）。这个特例让结果退化。

    \textbf{换一种设置}：让 $\ell$ 不过 $y$ 轴固定截距，改设 $\ell: y = kx + m$ 且 $m$ 可变，求 $|OA||OB|$ 最大时 $m, k$。完整讨论留给读者，以上只演示了弦长公式的机械流程。
}

\rmkb{
    做解析几何大题时，\textbf{一般路线}是：
    \begin{enumerate}
        \item 设直线方程（分斜率存在/不存在两种讨论）；
        \item 联立代入曲线，得到关于 $x$ 的二次方程；
        \item 用\textbf{韦达定理}得 $x_1 + x_2, x_1 x_2$；
        \item 用\textbf{弦长公式}、\textbf{距离公式}或\textbf{向量内积}把题目要求的量写成 $x_1 + x_2, x_1 x_2$ 的函数；
        \item 代入并化简，求最值或解方程。
    \end{enumerate}
    这是解析几何"韦达流"的完整套路。
}
